% page 536 du fac-simile (image p-535 du scan)
% bloc 157.5 x 237.7 mm ; marge gauche 8.0 mm
\pgc{-12.700mm}{0.000mm}{
\l{\cel{3}528.}
\l{}
\l{siliquo. (\sou{bot.}) Frukto longa, bivalva. - L. siliqua. EFIS.}
\l{}
\l{silko. Materio filamentoza, briloza, quan filifas speco di}
\l{\cel{3}faleno di qua la nomo Latina esas\cel{1}\sou{bombyx}. - ER.}
\l{}
\l{silogismo.\cel{2}(\sou{logiko}). Rezono formala, qua pruvas ke propozi-}
\l{\cel{3}ciono (\sou{konkluzo}) esas exakta, t.e. ke la atributo (\sou{termino}}
\l{\cel{3}\sou{majora}) di ta propoziciono konvenas a lua subjekto (\sou{ter}-}
\l{\cel{3}\sou{mino minora}) per la termino meza qua, konvenante ad la ter-}
\l{\cel{3}mino majora, ed a la termino minora, su kombinas kun singla}
\l{\cel{3}de li por formacar du propozicioni (majoro e minoro) -}
\l{\cel{3}la \textquotedbl{}premisi\textquotedbl{} - di qui la konkluzo esas la konsequo. - DE-}
\l{\cel{3}FIRS.}
\l{}
\l{\sou{silueto}. Desegnuro en qua streko trasita cirkum la ombro de}
\l{\cel{3}figuro, de korpo, reprezentas olua profilo,ed en qua to}
\l{\cel{3}quon cirkondas la konturi esas nigrigita. - DEFIRS.}
\l{}
\l{\sou{siluro}. (\sou{zool.)}\cel{2}Genero de fishi teleostea, tipo di la familio}
\l{\cel{3}\textquotedbl{}siluridi\textquotedbl{}, qui vivas en Europa ed en la regioni Aziana}
\l{\cel{3}di qui la klimato esas moderata - tre granda, kun kapo}
\l{\cel{3}plata, pelo glata e mola. - L.\cel{1}\sou{siluris glanis.}\cel{1}- EFIS.}
\l{}
\l{\sou{siluriano}. (\sou{geol.}) Nomo di la strati di tereno quin onu ren-}
\l{\cel{3}kontras sur la strati di olda greso reda. - DEFS.}
\l{}
\l{\sou{silvo}. Bosko en qua onu lasis la arbori atingar sua kresko-}
\l{\cel{3}punto maxima. - IL.}
\l{}
\l{\sou{silvano}. (\sou{mitol.}) Deo subalterna, deajo di la boski. - F.}
\l{}
\l{\sou{silvio}. (\sou{zool.}) Nomo vulgara di grupo de paseri, di qui la}
\l{\cel{3}plumi esas falvatra, e la kanto agreabla. - L.\cel{1}\sou{silvia}. -}
\l{}
\l{}
\l{\sou{simaro}. I. (en\cel{1}\sou{la epoki olima}). Longa robo di mulero. - II.}
\l{\cel{3}Sutano quan ula judiciisti weras sub sua robo oficala.}
\l{\cel{3}EFI.}
\l{}
\l{\sou{simbioso.}\cel{2}(\sou{biol.}) Asociitaro ek du o plura organismi, qui}
\l{\cel{3}interdiferas, kom speci, quan li profitas samatempe. (Kom}
\l{\cel{3}ex.: la asociitaro di algo kun fungo por formacar likeno).}
\l{\cel{3}DEFIRS.}
\l{}
\l{\sou{simbolo}. I. Objekto perceptebla, egardata kom la signo qua}
\l{\cel{3}reprezentas kozo qua ne esas perceptebla, per ula analoga-}
\l{\cel{3}jo qua frapis la imagino-fakultato. - II. (\sou{kemio}) Litero}
\l{\cel{3}vorto-komencala, adoptita por determinar korpo simpla en}
\l{\cel{3}la formuli kemiala. - III. (\sou{religio})\cel{2}Formulo konsakrita,}
\l{\cel{3}ek la kredendaji di religio. - DEFIRS.}
\l{}
\l{\sou{simboliko}. Cienco qua expozas ed interpretas la simboli qui}
\l{\cel{3}koncernas ta o ca religio o populo. -}
}
